\documentclass[11pt,a4paper]{article}
\usepackage[margin=25mm]{geometry}
\usepackage[T1]{fontenc}
\usepackage[utf8]{inputenc}
\usepackage{lmodern}
\usepackage{microtype}
\usepackage{amsmath,amssymb}
\usepackage{booktabs}
\usepackage{array}
\usepackage{longtable}
\usepackage{hyperref}
\usepackage{enumitem}
\usepackage{xcolor}
\usepackage{xurl}
\hypersetup{colorlinks=true,linkcolor=blue,citecolor=blue,urlcolor=blue}
\setlist{nosep,leftmargin=*}
\newcommand{\claimcount}[1]{\texttt{#1}}
\title{\Large A10 Claim-Bounded Black-Hole Interior Re-Exploration Archive\\[0.35em]
\normalsize Synthetic Diagnostic Candidate Mapping and External Comparator Handoff\\
\normalsize under the A10 Core Methodology}
\author{Keiji Yoshimura}
\date{May 28, 2026\\[0.5em]
\small GitHub Repository URL: https://github.com/yokken0907/a10-bh-claim-bounded-diagnostic-archive}
\begin{document}
\maketitle

\begin{abstract}
This manuscript reports a claim-bounded archive of a black-hole interior re-exploration route conducted under the A10 Core methodology. The work does not claim a black-hole interior solution, quantum extremal surface detection, singularity resolution, evidence for physics beyond general relativity, real stress-energy validation, numerical-relativity validation, or accepted physical candidates. Instead, the contribution is methodological and archival: previously over-pruned exploratory branches were reintroduced only as sandbox diagnostics, tested against negative controls and failure-boundary audits, and consolidated into an external-comparator handoff package. The locked route produced 216 sandbox diagnostic candidate events and two stable morphology clusters, while preserving zero physical claims and zero accepted physical candidates. Two important failures were detected and preserved: nondecaying-tail false accepts in an action/stress sandbox audit and collapse of paired morphology under an inappropriate global width gate. Both were remediated and retained as failure-boundary evidence. The final archive denies physical-entry authorization because 16 required external assets remain missing, including numerical-relativity comparators, gauge metadata, model-defined QFT/QES inputs, stress-energy/source consistency comparators, and theorem/no-go proof targets. The result is therefore best understood as a reproducible, claim-bounded diagnostic archive and specialist handoff brief, not as a report of new black-hole physics.
\end{abstract}

\section{Claim boundary and scope}

This paper is deliberately narrow. It reports a sandbox diagnostic archive, not a discovery of black-hole interior physics. The claim boundary is central to the result rather than a disclaimer appended after the fact. General relativity and black-hole solutions have a long, rigorous history beginning with the Einstein field equations and the Schwarzschild solution, with the Kruskal extension clarifying the maximal analytic extension of the Schwarzschild geometry \cite{einstein1915,schwarzschild1916,kruskal1960}. Numerical and semi-classical approaches to black-hole physics require careful treatment of gauge, boundary conditions, constraint residuals, and high-fidelity comparators \cite{gourgoulhon2007,pretorius2005,campanelli2006,baker2006}. The present archive does not meet those requirements for physical-entry claims.

The allowed claim is the following: a claim-bounded sandbox re-exploration pipeline was reconstructed under A10 Core v0.2.1; it produced a reduced diagnostic candidate map and an external-comparator handoff checklist while preserving zero physical claims and zero accepted physical candidates. The disallowed claims are equally important. This archive does not establish a black-hole interior solution, a quantum extremal surface (QES) detection, singularity resolution, beyond-GR evidence, real stress-energy validation, an Einstein-equation solution, a numerical-relativity validation, a semi-classical gravity solution, or an accepted physical candidate.

\section{Background and motivation}

The route was motivated by a change in A10 operating discipline. Earlier A10 black-hole exploration treated dangerous public claims and dangerous exploratory hypotheses as closely coupled. This was safe, but it may have over-pruned branches that could have been retained as non-physical diagnostics. A10 Core v0.2.1 separates entrance and exit conditions: speculative hypotheses may enter a sandbox route, but only evidence-bounded claims may exit it \cite{yoshimura2026a10core}. In this project, that distinction allowed entropy-extremal, action/stress, marching-solver, bridge-model, and handoff-comparator branches to be reopened without authorizing physical claims.

The phrase ``black-hole interior re-exploration'' should therefore be read as a route name, not as a claim that the black-hole interior has been physically accessed. Entanglement and extremal-surface ideas are well-established in holographic settings, from Ryu--Takayanagi surfaces to quantum extremal surfaces \cite{ryu2006,engelhardt2015}. However, this archive does not compute a model-defined QES in a controlled semi-classical spacetime. It only reopens a reduced entropy-extremal diagnostic branch as a sandbox component, then freezes a list of external assets that would be required before any physical interpretation could be attempted.

\section{Methodological design}

The route uses four principles.

\begin{enumerate}
  \item \textbf{A10-free baselines first.} Flat, Schwarzschild, Kruskal, and Vaidya-surrogate checks are replayed before A10-specific hypotheses are introduced.
  \item \textbf{Sandbox before physical entry.} Entropy-extremal, action/stress, bridge-model, and morphology branches are treated as reduced diagnostics, not physical candidates.
  \item \textbf{Failure boundaries are evidence.} Failed runs are archived and used to strengthen guards rather than erased.
  \item \textbf{External comparators are mandatory.} The route is not allowed to claim physical entry without real external assets and high-fidelity validation.
\end{enumerate}

This design follows the A10 pattern of bold hypothesis entry, strict process control, and bounded public claims. The key operational distinction is that a synthetic diagnostic event is not a physical candidate. It may guide future expert comparison, but it cannot by itself justify a physical statement.

\section{Route summary}

Table~\ref{tab:route} summarizes the major route blocks. The full route ledger is included in the repository as \texttt{data/R16\_route\_stage\_ledger.csv}.

\begin{table}[htbp]
\centering
\small
\caption{Condensed route summary. PASS does not imply physical-entry authorization.}
\label{tab:route}
\begin{tabular}{p{0.18\linewidth}p{0.34\linewidth}p{0.35\linewidth}}
\toprule
Block & Purpose & Locked interpretation \\
\midrule
R0--R1B & Re-exploration seed, A10-free baselines, grid and perturbation audits & Known-solution scaffolds established without physical claims. \\
R2--R2B & Entropy-extremal diagnostic revisit and negative controls & QES-like branch reopened only as reduced diagnostic. \\
R3--R3C & Action/stress sandbox and failure-boundary remediation & Nondecaying-tail false accepts detected and remediated. \\
R4--R4B & Marching-solver and randomized holdout checks & Manufactured solver stability verified in reduced setting. \\
R5--R7 & Candidate-map generation, adversarial controls, morphology clustering, consolidation & 216 diagnostic events and two stable morphology clusters locked. \\
R8--R12 & Pre-physical bridge, reduced residual, gauge, stress-proxy, and boundary sandbox checks & Bridge components specified but physical entry denied. \\
R13--R15P & External-comparator readiness, adapter/schema checks, asset binder, archive pause & 16 required external assets identified; all missing. \\
R16--R17 & Archive synthesis and repository preparation & Handoff-ready archive; not physical-entry-ready. \\
\bottomrule
\end{tabular}
\end{table}

\section{Locked diagnostic result}

The consolidated archive inherits the R16/R17 lock:

\begin{center}
\begin{tabular}{ll}
\toprule
Quantity & Locked value \\
\midrule
Route stage count & 30 \\
Diagnostic candidate event count & 216 \\
Stable morphology cluster count & 2 \\
Required external asset count & 16 \\
Bound required external asset count & 0 \\
Missing required external asset count & 16 \\
External comparator execution count & 0 \\
High-fidelity validation completed count & 0 \\
Physical-entry authorization & False \\
Physical claim count & 0 \\
Accepted physical candidate count & 0 \\
\bottomrule
\end{tabular}
\end{center}

The two sandbox morphology classes are \texttt{compact\_single\_like} and \texttt{paired\_compact\_like}. These names describe diagnostic shapes in the reduced archive. They do not name physical black-hole structures. Their only current function is to specify what future external comparators would need to test.

\section{Failure-boundary evidence}

Two failures are especially important because they prevent the archive from becoming a confirmation-only exercise.

\paragraph{R3B nondecaying-tail false accepts.}
The action/stress holdout audit initially accepted nondecaying-tail artifacts. This failure showed that the sandbox could mistake persistent tails for admissible structure. The route was therefore remediated through a nondecaying-tail guard and an integrated rerun. The failure is preserved in the ledger as \texttt{FAIL\textrightarrow REMEDIATED}, not hidden.

\paragraph{R6 paired-morphology collapse.}
A later synthetic-evolution coupling preflight collapsed the paired morphology class by applying an inappropriate global width gate. The issue was not a false physical discovery but a classification failure: a two-component morphology was being judged by a single-component compactness rule. The remediation introduced morphology-specific handling of pair separation and component width.

These failures support the archive's methodological value. They show where the route broke, why it broke, and how the guard logic was tightened without relaxing the claim boundary.

\section{External-comparator handoff}

The final blocker is not a vague lack of confidence. It is a concrete asset gap. The route identifies five external-comparator families:

\begin{enumerate}
  \item known-solution reference adapters;
  \item numerical-relativity metric comparators;
  \item semi-classical QFT/QES comparators;
  \item stress-energy/source consistency comparators;
  \item singularity-consistency theorem or no-go handoff routes.
\end{enumerate}

The repository lists 16 required assets across these families. They include reference datasets for flat, Schwarzschild, Kruskal, and Vaidya cases; numerical-relativity metric outputs with gauge metadata; model-defined QFT/QES calculation inputs; stress-energy/source consistency definitions; and theorem/no-go target statements. Since none of these assets are currently bound, the archive locks physical-entry authorization to false.

This is consistent with numerical-relativity practice: successful black-hole computation depends not only on code execution but also on formulation, gauge, constraints, boundary data, and validation against trusted comparators \cite{gourgoulhon2007,pretorius2005,campanelli2006,baker2006}. It is also consistent with the QES literature, where the object being extremized and the entropy functional must be model-defined before a QES claim can be meaningful \cite{ryu2006,engelhardt2015,bousso2022}.

\section{Discussion}

The route is best interpreted through a climbing analogy. The project did not climb the wall of black-hole interior physics. It reached the end of the hiking path, mapped the cliff edge, identified the required climbing equipment, and locked a safe turnaround point. In technical terms, the synthetic route produced a reproducible sandbox diagnostic map and a handoff specification, but not a physical solution.

This distinction matters because synthetic pipelines can easily produce plausible-looking patterns. Without strict negative controls, failure-boundary preservation, gauge and boundary audits, and external comparator requirements, such patterns can be mistaken for evidence. The present archive is designed to prevent that mistake. Its value is not that it makes a new physics claim, but that it keeps the exploratory route reproducible while refusing to authorize claims that its evidence cannot support.

\section{Limitations}

The limitations are structural.

\begin{itemize}
  \item No high-fidelity numerical-relativity validation has been executed.
  \item No real or trusted external metric dataset has been bound.
  \item No model-defined QFT/QES calculation has been performed.
  \item No real stress-energy tensor or semi-classical source has been validated.
  \item No theorem, no-go statement, or consistency proof for singularity behavior has been supplied.
  \item The diagnostic candidate map is synthetic and sandbox-limited.
\end{itemize}

Consequently, the appropriate next scientific step is not another synthetic phase. The appropriate next step is either archive/handoff publication or binding real external assets through a future \texttt{BH-R15C-bound-asset-smoke-adapter} route.

\section{Conclusion}

The A10-BH re-exploration archive has value as a claim-bounded diagnostic and handoff artifact, not as a black-hole physics discovery. It demonstrates that previously over-pruned exploratory branches can be reopened as sandbox diagnostics while preserving strict public claim boundaries. It consolidates 216 diagnostic candidate events into two stable morphology clusters, preserves and remediates two significant failure modes, and freezes the external assets required for any future physical-entry attempt. The archive therefore justifies GitHub repository publication and, with careful framing, may justify a preprint as a methodological case study in claim-bounded exploratory governance. It does not justify claims of black-hole interior discovery, QES detection, singularity resolution, or beyond-GR evidence.

\section*{Acknowledgements}

The author acknowledges the use of AI assistance, including ChatGPT, for iterative drafting, code-structure generation, claim-boundary review, and workflow organization. The author remains responsible for the final content, scope, and claims.

\begin{thebibliography}{99}

\bibitem{einstein1915}
A. Einstein, ``Die Feldgleichungen der Gravitation,'' \emph{Sitzungsberichte der Königlich Preußischen Akademie der Wissenschaften}, pp. 844--847, 1915.

\bibitem{schwarzschild1916}
K. Schwarzschild, ``Über das Gravitationsfeld eines Massenpunktes nach der Einsteinschen Theorie,'' \emph{Sitzungsberichte der Königlich Preußischen Akademie der Wissenschaften}, pp. 189--196, 1916.

\bibitem{kruskal1960}
M. D. Kruskal, ``Maximal Extension of Schwarzschild Metric,'' \emph{Physical Review}, vol. 119, no. 5, pp. 1743--1745, 1960. doi: \href{https://doi.org/10.1103/PhysRev.119.1743}{10.1103/PhysRev.119.1743}.

\bibitem{wald1984}
R. M. Wald, \emph{General Relativity}. University of Chicago Press, 1984.

\bibitem{gourgoulhon2007}
E. Gourgoulhon, ``3+1 Formalism and Bases of Numerical Relativity,'' arXiv:gr-qc/0703035, 2007.

\bibitem{pretorius2005}
F. Pretorius, ``Evolution of Binary Black-Hole Spacetimes,'' \emph{Physical Review Letters}, vol. 95, 121101, 2005. doi: \href{https://doi.org/10.1103/PhysRevLett.95.121101}{10.1103/PhysRevLett.95.121101}.

\bibitem{campanelli2006}
M. Campanelli, C. O. Lousto, P. Marronetti, and Y. Zlochower, ``Accurate Evolutions of Orbiting Black-Hole Binaries without Excision,'' \emph{Physical Review Letters}, vol. 96, 111101, 2006. doi: \href{https://doi.org/10.1103/PhysRevLett.96.111101}{10.1103/PhysRevLett.96.111101}.

\bibitem{baker2006}
J. G. Baker, J. Centrella, D.-I. Choi, M. Koppitz, and J. van Meter, ``Gravitational-Wave Extraction from an Inspiraling Configuration of Merging Black Holes,'' \emph{Physical Review Letters}, vol. 96, 111102, 2006. doi: \href{https://doi.org/10.1103/PhysRevLett.96.111102}{10.1103/PhysRevLett.96.111102}.

\bibitem{ryu2006}
S. Ryu and T. Takayanagi, ``Holographic Derivation of Entanglement Entropy from AdS/CFT,'' \emph{Physical Review Letters}, vol. 96, 181602, 2006. doi: \href{https://doi.org/10.1103/PhysRevLett.96.181602}{10.1103/PhysRevLett.96.181602}.

\bibitem{engelhardt2015}
N. Engelhardt and A. C. Wall, ``Quantum Extremal Surfaces: Holographic Entanglement Entropy beyond the Classical Regime,'' \emph{Journal of High Energy Physics}, 2015, 73, 2015. doi: \href{https://doi.org/10.1007/JHEP01(2015)073}{10.1007/JHEP01(2015)073}.

\bibitem{almheiri2020}
A. Almheiri, N. Engelhardt, D. Marolf, and H. Maxfield, ``The entropy of bulk quantum fields and the entanglement wedge of an evaporating black hole,'' \emph{Journal of High Energy Physics}, 2019, 63, 2019. doi: \href{https://doi.org/10.1007/JHEP12(2019)063}{10.1007/JHEP12(2019)063}.

\bibitem{bousso2022}
R. Bousso et al., ``Snowmass White Paper: Quantum Aspects of Black Holes and the Emergence of Spacetime,'' arXiv:2201.03096, 2022.

\bibitem{yoshimura2026a10core}
K. Yoshimura, \emph{A10 Core Foundational Methodology v0.2.1}, independent repository and methodology note, 2026. doi: \href{https://doi.org/10.5281/zenodo.20343694}{10.5281/zenodo.20343694}.

\end{thebibliography}

\appendix
\section{Repository files relevant to this manuscript}

The repository contains the following key supporting files. The manuscript should be read together with these files.

\begin{itemize}
  \item \path{docs/CLAIM_BOUNDARY.md}
  \item \path{docs/TECHNICAL_STATUS_LEDGER.md}
  \item \path{docs/SPECIALIST_HANDOFF_BRIEF.md}
  \item \path{data/R16_route_stage_ledger.csv}
  \item \path{data/R16_required_external_asset_checklist.csv}
  \item \path{static/upstream/R16_summary_lock.json}
\end{itemize}

\end{document}
